Work first in the uncapped Poisson comparison experiment.  Poisson thinning from \(M\sim\operatorname{Pois}(n/4)\) and the fair marks \(B_i\) gives, with \(m=n/8\), mutually independent counts
\[
 N'_{\jmath,x}\sim\operatorname{Pois}(mq_{\jmath,x}),\qquad
 N_{\jmath,x}\sim\operatorname{Pois}(mq_{\jmath,x})
\]
over all \((\jmath,x)\), with the pilot array independent of the evaluation array.  This uses \(\mathrm{ass:iid\mbox{-}sampling}\) and the Poisson splitting identity.  For a scalar \(N\sim\operatorname{Pois}(mq)\), the definition of \(U_h\) gives the exponential generating function
\[
 \sum_{h\geq0}U_h(N;z)\frac{w^h}{h!}=e^{-zw}(1+w/m)^N.
\]
Taking the expectation of one copy and of the product of copies at \(w,v\), respectively, gives
\[
 E e^{-zw}(1+w/m)^N=e^{(q-z)w},
\]
and
\[
 E\!\left[e^{-z(w+v)}(1+w/m)^N(1+v/m)^N\right]
 =\exp\{(q-z)(w+v)+(q/m)wv\}.
\]
Coefficient extraction therefore proves, conditionally on the pilot array,
\[
 E[U_h(N;z)]=(q-z)^h
\]
and
\[
 E[U_h(N;z)U_t(N;z)]
 =\sum_{\ell=0}^{h\wedge t}\binom h\ell\binom t\ell\ell!(q/m)^\ell(q-z)^{h+t-2\ell}.
\]

Set \(L=L_d\), \(\tau=L/m\), \(c_\jmath=N'_{\jmath,x}/m\), and
\[
 G_x=\bigcap_{\jmath\in\mathcal J}
 \bigl\{|c_\jmath-q_{\jmath,x}|\leq h_\jmath/4\bigr\},
 \qquad v_x=\sqrt{p_x\tau}+\tau.
\]
On \(G_x\), the definitions of \(\ell_\jmath,u_\jmath,b_\jmath,r_\jmath\) give \(q_x\in Q_x\) and \(h_\jmath/2\leq r_\jmath\leq h_\jmath\).  Solving the two quadratic inequalities relating \(c_\jmath\), \(q_{\jmath,x}\), and \(h_\jmath=H_0\{\sqrt{c_\jmath\tau}+\tau\}\) gives
\[
 c_\jmath\leq C(q_{\jmath,x}+\tau),\qquad
 q_{\jmath,x}\leq C(c_\jmath+\tau).
\]
Consequently,
\[
 \sum_\jmath r_\jmath\leq Cv_x,
 \qquad
 \frac{q_{\jmath,x}}{mr_\jmath^2}\leq\frac C L.                 \tag{4}
\]
These conclusions include \(q_{\jmath,x}=0\), because the corresponding Poisson count then vanishes almost surely.

For the complement of \(G_x\), let \(W\sim\operatorname{Pois}(\lambda)\).  Exponential Markov applied to
\[
 E e^{t(W-\lambda)}=\exp\{\lambda(e^t-1-t)\}
\]
gives, for \(z\geq1\),
\[
 \Pr\{|W-\lambda|>\sqrt{2\lambda z}+2z\}\leq2e^{-z}.            \tag{5}
\]
Choose the single numerical constant \(H_0\) large enough that the event complementary to that in (5), with \(z=100L\), implies
\[
 |W/m-q|\leq\frac{H_0}{4}
 \left\{\sqrt{(W/m)L/m}+L/m\right\}.
\]
Indeed, if \(\lambda\geq CL\), (5) implies \(W\geq\lambda/2\), and the square-root term dominates; if \(\lambda<CL\), the additive \(H_0L/m\) term dominates.  Integrating (5) above \(100L\), and using \(W/m\leq q+|W/m-q|\), proves for \(t=1,2,4\)
\[
 E\!\left[
 \left\{|W/m-q|+\sqrt{(W/m)\tau}+\tau\right\}^t
 \mathbf1_{\mathrm{bad}}\right]
 \leq C e^{-40L}(\sqrt{q\tau}+\tau)^t.                            \tag{6}
\]
Independence across the four coordinates, a union bound, and Cauchy--Schwarz with the fourth moment in (6) yield
\[
 E\!\left[
 \left\{\lVert b_x-q_x\rVert_1+\sum_\jmath r_\jmath\right\}^t
 \mathbf1_{G_x^c}\right]
 \leq C e^{-20L}v_x^t,
 \qquad t=1,2,                                                       \tag{7}
\]
whereas (4), together with the corresponding unconditional second moments, gives
\[
 E\!\left[\left(\sum_\jmath r_\jmath\right)^2\mathbf1_{G_x}\right]
 \leq Cv_x^2.                                                        \tag{8}
\]
Thus pilot failure retains the local \(p_x\)-scale.

Apply \(\mathrm{lem:simultaneous\mbox{-}jackson\mbox{-}certificate}\) on \(G_x\).  If \(\ell_\jmath=0\), then \(c_\jmath\leq h_\jmath\), so \(c_\jmath+h_\jmath\leq C\tau\) and
\[
 \sqrt{(q_{\jmath,x}-\ell_\jmath)(u_\jmath-q_{\jmath,x})}
 \leq C\sqrt{q_{\jmath,x}\tau},
 \qquad r_\jmath\leq C\tau.
\]
If \(\ell_\jmath>0\), then \(c_\jmath>h_\jmath\), so \(q_{\jmath,x}\geq3c_\jmath/4\) and \(r_\jmath\leq C\sqrt{q_{\jmath,x}\tau}\).  Since \(\sum_\jmath\sqrt{q_{\jmath,x}}\leq2\sqrt{p_x}\), the vertex-adaptive Jackson bound gives
\[
 |P_{x,K_d}(q_x)-\bar f_\epsilon(q_x)|
 \leq C_\epsilon\left\{
 \frac{\sqrt{p_x\tau}}{K_d}+\frac{\tau}{K_d^2}\right\}.          \tag{9}
\]

Conditional on the pilot, factorial unbiasedness gives
\(E(Z_x\mid N')=P_{x,K_d}(q_x)\).  From the second factorial identity and (4), whenever \(|q-q_0|\leq r\),
\[
 \frac{E\{U_h(N;q_0)^2\}}{r^{2h}}
 \leq\sum_{\ell=0}^h\binom h\ell^2\ell!(C/L)^\ell
 \leq\exp(Ch^2/L).                                                   \tag{10}
\]
Combining (10), coordinate independence, Minkowski's inequality, and the coefficient envelope in \(\mathrm{lem:simultaneous\mbox{-}jackson\mbox{-}certificate}\), and writing
\[
 S_x=(1+\epsilon^{-1})\sum_\jmath r_\jmath,
\]
gives
\[
 E\{(Z_x-\bar f_\epsilon(b_x))^2\mid N'\}
 \leq S_x^2\exp(C_1K_d+C_2K_d^2/L).                                \tag{11}
\]
Choose \(\kappa>0\) sufficiently small and then \(D_0\) sufficiently large.  Since \(K_d=\max\{2,\lfloor\kappa L\rfloor\}\), (11) is at most \(CS_x^2d^{1/16}\) for \(d\geq D_0\); increasing \(C\) covers the remaining finite values.

Write \(W_x=Z_x-\bar f_\epsilon(b_x)\), \(t_x=d^{1/4}S_x\).  The elementary projection inequalities
\[
 |\Pi_{[-t,t]}(w)-w|\leq w^2/t,
 \qquad |\Pi_{[-t,t]}(w)|\leq|w|                                  \tag{12}
\]
hold for every \(t>0\), with the case \(t=0\) understood by continuity; in that case every radius and every relevant cell mass is zero.  On \(G_x\), equations (9)--(12), global Lipschitzness from \(\mathrm{prop:identification\mbox{-}and\mbox{-}extension}\), and \(d^{-3/16}\leq C/L^2\) give
\[
 |E(T_x^{\mathrm{JF}}\mid N')-\bar f_\epsilon(q_x)|
 \leq C_\epsilon\left\{
 \sqrt{\frac{p_x}{mL}}+\frac1{mL}\right\},                       \tag{13}
\]
and
\[
 E\{(T_x^{\mathrm{JF}}-\bar f_\epsilon(q_x))^2\mid N'\}
 \leq CS_x^2d^{1/16}.                                               \tag{14}
\]
On \(G_x^c\), clipping and global Lipschitzness give the deterministic bound
\[
 |T_x^{\mathrm{JF}}-\bar f_\epsilon(q_x)|
 \leq(1+\epsilon^{-1})\lVert b_x-q_x\rVert_1+d^{1/4}S_x.
\]
Its first two moments are controlled by (7); since \(e^{-20L}d^{1/2}\leq Cd^{1/16}\), those moments fit, respectively, into the right side of (13) and the expectation of (14).  Integrating (13)--(14), using (8), and using
\[
 v_x^2=(\sqrt{p_xL/m}+L/m)^2
 \leq2\{p_xL/m+L^2/m^2\}
\]
proves
\[
 |E T_x^{\mathrm{JF}}-\bar f_\epsilon(q_x)|
 \leq C_\epsilon\left\{\sqrt{\frac{p_x}{mL_d}}+\frac1{mL_d}\right\},
\]
and
\[
 \operatorname{Var}(T_x^{\mathrm{JF}})
 \leq C_\epsilon d^{1/16}
 \left\{\frac{p_xL_d}{m}+\frac{L_d^2}{m^2}\right\}.
\]
This proves the stated control, including pilot failure, null coordinates, boundary means, and treatment-effect ties.